\paragraph{Alternation with three owners.}
One can use just two owners for resets along the marked path. At a
disjunction, compare against the union of two $b$-coalitions. A helper
that wins on this union need not win on either part. Changing one
part to $c$ makes the otherwise untested case visible as a winner
that must be erased later on the path.

\begin{theorem}[Three-owner alternating bounds]
\label{core:thm:three-owner-alternating-bounds}
For every fixed $K\ge3$, supplying one positive under source bounds
given by their canonical prime lists is EXPTIME-hard under
polynomial-time many-one reductions, even at bounded prime width.
The construction below is feasible exactly when the start of the
state game wins. If its winning height is $H$, the least number of
owner slots is between $2H-1$ and $2H+2$.
\end{theorem}
\begin{proof}
Use owners $a,b,c$. The main spine uses $a,b$, and the requested
source is $X$ with all rails at $c$. Let $s$ denote a base state.
An $a$-address $u=(s,\epsilon,p)$ includes a phase bit $\epsilon$ and
either no previous address or an address $p$ in the base state bits.
Use the following disjoint groups of rail pairs.
\begin{itemize}
\item $A_u$ selects the $a$-address on the primary rails.
\item $L_i(s,\epsilon)$ for $i=0,1$ selects the base state and phase
on the two left groups. Use both at conjunctions and only $L_0$ at
disjunctions.
\item $R_\epsilon(s)$ and $Z_\epsilon(s)$ select the base state on
two groups for that phase. Use these at disjunctions.
\end{itemize}
Write a bar for the opposite rails of an address within its group.
The test for an $a$-reset at $u$ posts the coalition
$\widehat A_u=A_u\cup\overline R_{1-\epsilon}(p)$ at $a$, omitting
the second set when $p$ is absent. On taking a left move $i$ from
$(s,\epsilon,p)$ the successor address is
$(s_i,1-\epsilon,\varnothing)$. On taking a right move from a
disjunction it is $(s_1,1-\epsilon,s)$. Write these addresses as
$u_i^{\rm L}$ and $u^{\rm R}$. Use the seed loss posting the initial
$A$-address at $a$, with phase zero and no previous address.

At a rail head in $A$ or $L_i$ require at least one $c$-posting in
each pair of its group. When there is exactly one per pair, the
other rails give the address including the head. At an $A$-head,
require $L_0$ or $L_1$ fully at $b$ at a conjunction. At a
disjunction require $L_0$, $R_\epsilon(s)$ and the entire group
of $Z_\epsilon$ rails at $b$. A rejecting state loses with no
further posting and an accepting state has no loss. At an
$L_i$-head require $A_{u_i^{\rm L}}$ at $a$ on transitions that use
this left group. The successor address does not depend on $p$.

At an $R_\epsilon$-head impose a loss when all $Z_\epsilon$ rails
are $c$. In addition, require a loss when each $R_\epsilon$ pair
has an $a$-posting and the unique remaining address, when unique,
is a disjunction whose right successor $A_{u^{\rm R}}$ is posted
at $a$. At a $Z_\epsilon$-head impose a loss when all
$R_\epsilon$ rails are $c$, and the same right-successor loss
using $b$-postings to specify the remaining $Z_\epsilon$ address.
Invalid or inapplicable addresses have no losses under the address
rules. Extend these rules antitonically at doubled postings, keeping
their at-least-one-per-pair tests. The two unconditional tests on
the other groups need no address.

For necessity, linearize an occurrence of $X$ as in
\cref{core:thm:alternating-rail-transitions}. Let $L_a(u)$ be the
last attachment winning on $\widehat A_u$. The seed forces
$L_a$ at the start. The positive excludes every $c$-attachment
on the marked spine. A comparison to omit $r$ gives $X$ the
full complement of $r$, so the spine inherits the omission.
For the comparisons omitting $a$ or $b$, fill in unused rail values
with $\{c\}$. At any comparison omitting $a$, the extra $a$-sources
from $\overline R_{1-\epsilon}(p)$ are suppressed by the condition
on the previous phase's $Z$ group.

At a rejection, compare with just $\widehat A_u$ posted at $a$.
There is no foreign reset available to remove $a$, so $L_a(u)$
cannot exist. At a conjunction use the two comparisons posting
$\widehat A_u$ and $L_i(s,\epsilon)$ at their owners and omitting
$a$. Each forces a later $b$-reset on its stated $L_i$. The
successor comparisons keep those $b$-flags, post the corresponding
$A_{u_i^{\rm L}}$ at $a$, and omit $b$, forcing still later
$a$-resets. In particular both successors are required.

At a disjunction put
\[
  B=L_0(s,\epsilon)\cup\overline Z_\epsilon(s),\qquad
  C=R_\epsilon(s)\cup Z_\epsilon(s).
\]
The comparison with $\widehat A_u$ at $a$ and $B\cup C$ at $b$
omits $a$ and forces a $b$-attachment with helper $h$ at a position
strictly after $L_a(u)$. Its $b$-test wins on $B\cup C$.
There are two successor comparisons.
\begin{itemize}
\item Post $B$ at $b$ and $A_{u_0^{\rm L}}$ at $a$ and omit
$b$. The sources in $\overline Z_\epsilon(s)$ are suppressed
by the condition on the $R_\epsilon$ group. A $b$-reset on $B$
after $L_a(u)$ therefore forces the left successor later still.
\item Post $B$ at $b$, $C$ at $c$ and $\widehat A_{u^{\rm R}}$
at $a$. Give every other input the full complement of $c$.
Exactly the sources in $C$ can supply $c$, and the address
losses omit them. The $a$-flags here are precisely
$\widehat A_{u^{\rm R}}$.
\end{itemize}
At the second comparison, immediately after $L_a(u)$ the value
cannot be $\{b\}$. Its helper there cannot be $\{b\}$, since the
winning $a$-coalition $\widehat A_u$ is disjoint from $B$ and
distinct owners' winning coalitions must intersect. A foreign
attachment can give $\{b\}$ only when both operands do.
If the helper $h$ in this comparison omits $c$, replacing its
$C$ inputs by $b$ preserves its output by omission independence.
That replacement makes $h$ a $b$-atom, so the original $h$ was a
$b$-atom too. If $h$ contains $c$ and the incoming value is
$\{b\}$, some $b$-helper has already reset on $B$ after $L_a(u)$.
In either case the first successor comparison yields the left
successor. In the remaining case the attachment of $h$ transmits
$c$. Removing it requires a later $b$- or $a$-reset, giving the
left or right successor, respectively. Thus a disjunction requires
one strictly later successor. Finite position chains give winning
trees, with at least $2H-1$ positions on a longest branch.

For sufficiency choose decreasing-height moves through the winning
base states. At each height attach the $c$-folds on $A_u$ for the
states of that height, including all their phases and previous
addresses. The extra rails in $\widehat A_u$ are used only in the
comparisons. In the following $b$-slot attach the chosen
$c$-folds on $L_i(s,\epsilon)$ for left moves, and for a right move
attach the helper
\[
   \left(\mathop{\star_a}\limits_{v\in R_\epsilon(s)} v\right)
       \star_c
   \left(\mathop{\star_b}\limits_{v\in Z_\epsilon(s)} v\right).
\]
On the disjunction's joint $b$-posting the right helper gives
$\{b\}$. At an $A$-loss the chosen competitor resets later. At
an $L_i$-loss its successor at lower height resets later.
On a loss of an $R$-head with all $Z$-rails at $c$, or of a
$Z$-head with all $R$-rails at $c$, the helper for that phase
gives $\{c\}$. Such a loss is visible only to other owners.
At a right-address loss only a helper with the remaining address
can transmit the head's mark. Every other helper has an absorbing
$a$-input in its $R$ part, or an absorbing $b$-input in its $Z$
part, respectively. Either absorption removes a mark from that
part on the observer facet where the loss is visible. A doubled
pair suppresses the mark in every helper of that group. When
the addressed right helper is used the required successor
$a$-reset is later. These are all the rail losses. The seed
absorbs $X$ at its loss, and the all-$c$ probe retains $X$.
There are at most $2H-1$ main slots and three more common slots
for the helper folds, using unary padding.

For the complexity bound use the local machine rules in
\cref{core:cor:exptime-bounded-source}. The retained base state,
the phase bit and the previous address are written by local rules
with bounded antecedents. Each address rule above uses only
one posting color for each input at a given head. Expand the
conjunction rules at the $A$-heads and the successor bit rules
as in that corollary. At an $R$- or $Z$-head the added loss is a
conjunction of unit clauses using inputs disjoint from the
address rule. Distributing this disjunction adds at most one
posting per clause and only polynomially many clauses overall.
Non-addressed phases and invalid codes can be excluded by the
same local clauses. These formulas give admissible extensions
on doubled pairs by the absorbing tests just proved. All
postings within a head still use a fixed color per input, so
facet minimization gives the complete canonical lists.
\end{proof}
